\documentclass[journal]{IEEEtran}

\usepackage{cite}
\usepackage{amsmath,amssymb,amsfonts}
\usepackage{bm}
\usepackage{graphicx}
\usepackage{adjustbox}
\usepackage{booktabs}
\usepackage{multirow}
\usepackage{array}
\usepackage{hyperref}
\usepackage{textcomp}

\hypersetup{hidelinks}

\def\BibTeX{{\rm B\kern-.05em{\sc i\kern-.025em b}\kern-.08em
    T\kern-.1667em\lower.7ex\hbox{E}\kern-.125emX}}
\markboth{IEEE Transactions on Medical Imaging}%
{Anonymous Author(s): SurgAlign}

\title{SurgAlign: Learning from Temporally Misaligned Weak Surgical Video-Text Pairs}
\author{Anonymous Author(s)}

\begin{document}

\maketitle

\begin{abstract}
Large-scale video-language pretraining often relies on weakly aligned video-text pairs whose timestamps are only approximate. In surgical video, this mismatch is especially severe: a caption may be correct for a coarse interval even though the frames that visually justify it occupy only a brief moment within repetitive procedural context. Standard CLIP-style training pools sampled frames from the weak interval into one positive clip, which mixes supportive and non-supportive content. We propose \textbf{SurgAlign}, a training-time method that treats each weak interval as a noisy temporal anchor, expands it into a small local window, and uses the paired text to softly weight candidate frames while keeping inference unchanged. Under a matched surgical-domain protocol with fixed data, encoders, and optimization, SurgAlign improves a matched baseline on several zero-shot benchmarks, with the largest gain on the evidence-sensitive GraSP Step task ($+5.43$ Acc / $+5.27$ F1). The results show that under weak temporal supervision, better construction of positive video-text pairs can materially improve zero-shot surgical understanding.
\end{abstract}

\begin{IEEEkeywords}
Contrastive learning, surgical video understanding, temporal misalignment, video-language pretraining, weak supervision.
\end{IEEEkeywords}

\section{Introduction}

Understanding surgical videos is fundamental to workflow recognition, context-aware assistance, retrospective review, quality assessment, and surgical education. Yet the labels required by supervised systems, including phases, steps, actions, tools, and triplets, remain expensive to annotate at scale. This bottleneck has made surgical video-language pretraining (VLP) increasingly attractive: instead of learning each downstream task from scratch, a model can first learn transferable surgical semantics from weakly paired video-text data and then support zero-shot or low-supervision evaluation.

Recent surgical VLP systems have shown that this strategy works. Lecture videos, narrated procedures, and curated corpora such as SurgLaVi provide scalable supervision and have enabled steady progress in surgical representation learning~\cite{lecturevideo2023,peskavlp2024,surglavi2025}. However, most existing pipelines still inherit a standard CLIP-style assumption: once a text span is paired with a video interval, the sampled clip is treated as a reliable positive visual unit for the paired text. Multiple frames are sampled from the annotated interval, pooled into one visual representation, and aligned to the text through a symmetric contrastive loss.

This assumption is convenient, but in surgical data it is often inaccurate. Our text supervision is typically not a manually curated caption for a cleanly trimmed clip. Instead, it is derived from narrations, transcripts, or weak temporal descriptions attached to a broader interval. In such settings, the textual evidence can be temporally offset from the most informative visual moment: a surgeon may describe the next action before it happens, summarize a broader sub-procedure after the key manipulation has passed, or narrate a local event whose discriminative visual evidence occupies only a short subsegment of the annotated interval. Similar weak synchronization has been reported in instructional and narrated videos beyond surgery~\cite{shen2021dwsa,jin2023ground_articles,xu2023procedure_aware_repr}. The result is a practical mismatch: a video-text pair may be globally valid while still containing only a small subset of frames that truly support the text.

This mismatch matters precisely because modern surgical benchmarks are not purely coarse-grained. Phase recognition can benefit from broad contextual cues, but step, action, tool, and triplet prediction rely much more heavily on localized visual evidence. If weak supervision is always consumed through uniform interval-level pooling, the representation may drift toward broad procedural context at the expense of the local cues required for fine-grained transfer. In other words, the problem is not simply noisy annotation; it is the training assumption that all frames inside a weakly paired interval are equally useful positives.

We argue that this problem should be addressed directly as \emph{temporal misalignment inside weak video-text pairs}. This view is consistent with ideas from multiple-instance learning and weak temporal localization, where a coarse positive unit is better treated as a partially positive bag than as a uniformly valid set of instances~\cite{dietterich1997mil,ilse2018attentionmil,li2023boostingwtal,hu2024mmplateau,li2024fps}. It is also aligned with recent video-language work showing that fixed temporal pooling often mixes relevant and irrelevant evidence, and that text-guided aggregation can recover more informative local correspondence~\cite{gorti2022xpool,li2023prost,liu2023stan}. In surgery, this observation becomes particularly important because temporal transitions are subtle, workflow descriptions are semantically dense, and narration-style supervision is rarely perfectly synchronized with visual events.

Based on this perspective, we present \textbf{SurgAlign}, a misalignment-aware training framework for surgical VLP. SurgAlign treats each weak interval as a noisy temporal anchor rather than a trustworthy clip boundary. During training, it expands the interval into a bounded temporal neighborhood, samples candidate frames from the expanded window, and uses the paired text to softly reweight local visual evidence. The resulting selected view is trained jointly with a standard global branch, so the model retains a stable global alignment signal while also learning to recover text-relevant local evidence from weakly aligned intervals. Crucially, we use multi-level annotations to condition this evidence-selection process itself. Fine, mid, and coarse captions differ not only in temporal scope but also in internal semantic complexity: fine captions are often dominated by a single cue, whereas broader captions more often behave as compositional semantic structures. SurgAlign therefore uses hierarchical labels to jointly control how many latent semantic cues may be extracted from the text and how sharply the temporal evidence is concentrated. In the hierarchical setting, we additionally regularize selected views from same-video adjacent levels so that local evidence remains compatible across nested temporal scopes.

This design is deliberately lightweight. SurgAlign does not introduce a new test-time retrieval protocol, a new surgical foundation model, or a heavy query-conditioned matching module. All additional logic is used only during pretraining, and inference remains standard CLIP-style independent video and text encoding. This makes the method easy to position: our contribution is not a new backbone, but a better way to consume weak temporal supervision on top of existing backbones.

The paper is organized around two complementary empirical questions. First, can misalignment-aware temporal denoising improve surgical VLP when compared under a strong surgical-domain backbone? Second, does the same training rule remain useful when the initialization backbone changes? We therefore use a surgical-domain primary model for the main benchmark comparison, and separately provide a cross-backbone validation to show that the gain is not tied to a single initialization family.

In summary, this paper makes three contributions. First, it identifies temporal misalignment inside weakly paired surgical video-text samples as an under-modeled problem in current surgical VLP. Second, it proposes SurgAlign, a training framework that combines expanded temporal evidence with multi-level annotation-aware evidence selection, where hierarchy jointly conditions latent semantic cue extraction and temporal concentration while preserving standard CLIP-style inference. Third, it validates the method both in a primary surgical-domain setting and in a cross-backbone setting, supporting the claim that the contribution lies in the training strategy rather than in a particular backbone choice.


\section{Related Work}
\label{sec:related_work}

\paragraph{Supervision bottlenecks in surgical video understanding.}
Surgical video understanding has often been studied through supervised tasks such as phase recognition, workflow modeling with timestamp labels, and action-triplet recognition~\cite{tecno2020,lessismore2022,rendezvous2022}. These settings have driven substantial progress, but they also highlight the practical bottleneck of surgical AI: obtaining dense expert annotations for phases, steps, tools, actions, and targets is costly, center-specific, and difficult to scale across procedures. This limitation motivates representation learning paradigms that can exploit weak supervision and support more transferable or zero-shot evaluation.

\paragraph{Surgical video-language pretraining.}
Against this background, surgical video-language pretraining has become an important direction for scalable and transferable surgical representation learning. Early work such as SurgVLP established surgical lecture videos as a viable weak supervision source for CLIP-style pretraining~\cite{lecturevideo2023}. Later studies extended this line through hierarchy-aware supervision, procedure-aware augmentation, retrieval-based enhancement, and larger-scale curated corpora, including HecVL, PeskaVLP, OphCLIP, and SurgLaVi~\cite{hecvl2024,peskavlp2024,ophclip2024,surglavi2025}. These works collectively show that weak video-text supervision can improve zero-shot or low-supervision surgical understanding, especially when the textual side becomes larger, cleaner, and more hierarchically structured.

\paragraph{Alignment reliability under weak surgical supervision.}
At the same time, recent surgical VLP papers make clear that richer text does not automatically imply better alignment. SurgVLP notes semantic misalignment between clips and text in surgical lecture videos~\cite{lecturevideo2023}, SurgLaVi highlights noisy pairs caused by temporal misalignment and commentary drift~\cite{surglavi2025}, and PeskaVLP identifies textual information loss and cross-modal procedural alignment as key challenges~\cite{peskavlp2024}. However, most CLIP-style surgical VLP pipelines still optimize standard contrastive learning over selected clip-caption pairs, effectively treating the sampled interval as a single positive unit once the pair has been constructed. This leaves open a more specific question: how to denoise weak temporal positives inside an otherwise valid clip-caption pair.

\paragraph{Partial-positive learning and text-guided evidence selection.}
Outside the surgical domain, weak video-text learning has increasingly been treated as an alignment-quality problem rather than a simple matter of text availability. Multiple-instance learning provides a useful conceptual lens: a weakly labeled bag can be positive even when only a subset of instances truly supports the label~\cite{dietterich1997mil,ilse2018attentionmil}. In video understanding, weakly supervised temporal localization further shows that language can help suppress irrelevant snippets and foreground latent evidence under coarse supervision~\cite{li2023boostingwtal,hu2024mmplateau,li2024fps}. Text-conditioned aggregation methods such as X-Pool similarly show that fixed mean pooling can be replaced by relevance-aware weighting~\cite{gorti2022xpool}. Our window denoising branch is most closely aligned with this line of work: we treat the expanded interval as a bounded partially positive bag and use the paired caption to estimate latent evidence weights during training.

\paragraph{Hierarchy as a granularity prior.}
Hierarchical video-language datasets and models such as SurgLaVi, HiTeA, and HecVL show that temporal granularity should not be ignored during representation learning~\cite{surglavi2025,ye2023hitea,hecvl2024}. In our work, however, hierarchy is not the main supervision target. Instead, we use fine/mid/coarse labels only as a lightweight prior on expected evidence concentration during denoising: fine captions should induce sharper frame weighting, whereas broader captions should permit smoother evidence distributions.

\paragraph{Position of our work.}
Our work is closest to the surgical CLIP-style pretraining line, but differs in what it treats as the unresolved bottleneck. We do not primarily contribute a larger corpus, a new backbone, or a new hierarchical objective. Instead, we focus on how weak temporal supervision is consumed during training once a pair has already been constructed. Specifically, we argue that a globally valid weak pair can still contain temporally noisy positives, and we address that issue with train-time window denoising while preserving standard CLIP-style inference at test time.


\section{Method}
\label{sec:method}

Our goal is to improve how weak temporal supervision is used during surgical VLP without replacing the underlying CLIP-style model. The key assumption we relax is simple: a weakly paired interval should not be treated as a uniformly reliable positive clip. Instead, we treat it as a noisy temporal anchor whose supporting evidence may lie in only part of the interval and may be slightly shifted around the annotated boundary. SurgAlign addresses this issue through training-time temporal denoising.

\subsection{Problem Formulation}
\label{sec:method_setup}

Each training sample is written as
\begin{equation}
(V, [t_s, t_e], y, \ell),
\end{equation}
where $V$ is a surgical video, $[t_s, t_e]$ is the weak temporal interval, $y$ is the paired text, and $\ell \in \{f, m, c\}$ denotes the fine, mid, or coarse annotation level. In a standard CLIP-style baseline, $M$ frames are sampled from $[t_s, t_e]$, encoded by a visual backbone $E_v$, aggregated into one clip representation, and aligned with a text representation produced by a text encoder $E_t$. If $A_v(\cdot)$ is the temporal aggregation head and $P_v(\cdot), P_t(\cdot)$ are the projection heads, the baseline computes
\begin{equation}
v_{\mathrm{glob}} = \operatorname{norm}\!\left(P_v(A_v(E_v(c)))\right), \qquad
t = \operatorname{norm}\!\left(P_t(E_t(y))\right),
\end{equation}
and optimizes a symmetric contrastive loss.

This formulation is effective when the sampled clip is already well aligned with the text. Our setting is weaker: the pair may be globally correct while the truly supportive visual evidence occupies only a subset of frames. SurgAlign keeps the same backbone and contrastive formulation, but changes how the training clip is constructed and weighted.

\subsection{Overview}
\label{sec:method_overview}

The overall pipeline is illustrated in the framework diagram. Starting from a weak temporal interval, SurgAlign first expands the interval into a bounded neighborhood. It then samples candidate frames from this expanded window and extracts frame-level features. Instead of uniformly pooling these frames, SurgAlign computes text-conditioned frame scores, converts them into soft weights, and builds a selected local visual view. The model is trained with two complementary branches:
\begin{enumerate}[leftmargin=*]
    \item a global branch that preserves stable interval-level video-text alignment; and
    \item a selected branch that emphasizes text-relevant evidence inside the expanded temporal neighborhood.
\end{enumerate}
In hierarchical training, SurgAlign also encourages adjacent levels from the same video to produce compatible selected views. All of these mechanisms are training-only. Test-time inference remains standard CLIP-style video/text encoding.

\subsection{Expanded Temporal Evidence}
\label{sec:method_expand}

Weak narration timestamps often miss the most informative visual moment by a small temporal offset. We therefore replace the original interval with a bounded expanded interval
\begin{equation}
[t_s', t_e'] = \operatorname{Expand}([t_s, t_e]; \rho),
\end{equation}
where $\rho \geq 1$ is an expansion ratio and the interval is clipped to the valid video duration. We then sample a multi-frame candidate clip
\begin{equation}
c_{\mathrm{exp}} = \{x_1', \ldots, x_M'\}
\end{equation}
from $[t_s', t_e']$.

This step is deliberately modest. SurgAlign does not attempt full temporal grounding. The purpose of expansion is to give the model a small search space around a noisy timestamp so that the true supporting frames remain available during training even when the annotated boundary is slightly off.

\subsection{Multi-Level Annotation-Aware Evidence Selection}
\label{sec:method_selection}

Let the visual encoder produce frame tokens
\begin{equation}
\{h_1, \ldots, h_M\} = E_v(c_{\mathrm{exp}}),
\end{equation}
and let the text encoder produce token-level hidden states and a pooled global text representation
\begin{equation}
\{w_1, \ldots, w_L\}, \quad t = E_t(y).
\end{equation}
The global text embedding $t$ is still used by the contrastive objective. The selected branch, however, uses the token-level states to construct text-conditioned temporal evidence inside the expanded window. Our key design choice is to treat the multi-level annotation label $\ell$ as a prior on \emph{how} this evidence should be selected. Fine-, mid-, and coarse-level captions differ not only in temporal scope but also in internal semantic complexity. Fine captions are usually localized and dominated by one salient cue, whereas mid/coarse captions more often summarize multiple actions, instruments, or procedural states within one interval. We therefore treat mid/coarse captions as \emph{compositional semantic structures} rather than as longer monolithic queries.

We encode this prior through a level-conditioned evidence-selection rule. Given level $\ell$, the selected branch uses a pair of level-dependent parameters,
\begin{equation}
K_\ell, \qquad \tau_{\mathrm{sel}}^{(\ell)},
\end{equation}
where $K_\ell$ controls how many latent semantic cues may be extracted from the text and $\tau_{\mathrm{sel}}^{(\ell)}$ controls how sharply the temporal evidence is concentrated. Concretely, we use
\begin{equation}
K_\ell =
\begin{cases}
K_f, & \ell = f,\\
K_m, & \ell = m,\\
K_c, & \ell = c,
\end{cases}
\qquad
\tau_{\mathrm{sel}}^{(\ell)} =
\begin{cases}
\tau_f, & \ell = f,\\
\tau_m, & \ell = m,\\
\tau_c, & \ell = c,
\end{cases}
\end{equation}
with $K_f \leq K_m \leq K_c$ and $\tau_f < \tau_m < \tau_c$. This means that fine captions use fewer queries and sharper temporal selection, while broader captions are allowed to mine multiple latent semantic cues and distribute their mass over a wider temporal neighborhood.

To keep the training objective simple, we do not explicitly decompose the sentence into supervised sub-phrases. Instead, we introduce a small bank of learnable latent text-query slots
\begin{equation}
\{q_1, \ldots, q_{K_{\max}}\},
\end{equation}
which are shared across levels and softly attend to the token sequence to extract multiple latent semantic cues.

For the $k$-th latent query, we compute token attention
\begin{equation}
\beta_{kr} =
\operatorname{softmax}_{r}
\left(
\frac{q_k^\top w_r}{\sqrt{d}}
\right),
\end{equation}
and obtain a query-specific text cue
\begin{equation}
u_k = \sum_{r=1}^{L} \beta_{kr} w_r .
\end{equation}
These latent cues are then matched against the frame tokens after lightweight projections:
\begin{equation}
\tilde{u}_k = \operatorname{norm}(Q_t(u_k)), \qquad
\tilde{h}_j = \operatorname{norm}(Q_v(h_j)).
\end{equation}
Their similarity defines query-conditioned frame scores
\begin{equation}
s_{kj} = \tilde{u}_k^\top \tilde{h}_j .
\end{equation}

Given the annotation level $\ell \in \{f,m,c\}$, we activate only the first $K_\ell$ latent queries and normalize the scores over time with the corresponding selection temperature:
\begin{equation}
\alpha_{kj} =
\operatorname{softmax}_{j}
\left(
\frac{s_{kj}}{\tau_{\mathrm{sel}}^{(\ell)}}
\right), \qquad k = 1, \ldots, K_\ell .
\end{equation}
Each latent query therefore selects its own temporally weighted evidence from the expanded window. We aggregate the query-specific evidence as
\begin{equation}
z_k = \sum_{j=1}^{M} \alpha_{kj} h_j,
\end{equation}
and fuse them into a single selected visual view
\begin{equation}
z_{\mathrm{sel}} = \frac{1}{K_\ell}\sum_{k=1}^{K_\ell} z_k, \qquad
v_{\mathrm{sel}} = \operatorname{norm}(P_v(z_{\mathrm{sel}})).
\end{equation}

This selected branch should be interpreted as a denoised auxiliary view of the same weak pair, not as a new inference-time matching module. The global branch still provides stable clip-level supervision, while the selected branch teaches the model to recover text-relevant temporal evidence from imprecise positives. Importantly, multi-level annotation is not used to define a separate hierarchy-specific prediction task. It simply conditions the evidence-selection process through $(K_\ell, \tau_{\mathrm{sel}}^{(\ell)})$. The role of multiple latent queries is therefore not to enforce explicit sentence parsing, but to provide a lightweight evidence-mining mechanism when the caption expresses a compositional procedure description.

\subsection{Hierarchical Consistency Across Adjacent Levels}
\label{sec:method_consistency}

When fine, mid, and coarse samples from the same video are jointly observed during training, their selected views should remain compatible rather than drifting arbitrarily in the embedding space. We therefore construct mixed-level batches and, when triplet anchoring is enabled, include same-video nested triplets satisfying
\begin{equation}
[t_s^{f}, t_e^{f}] \subseteq [t_s^{m}, t_e^{m}] \subseteq [t_s^{c}, t_e^{c}].
\end{equation}
For the available triplets $\mathcal{T}$ in a batch, we apply an adjacent-level consistency loss
\begin{equation}
\mathcal{L}_{\mathrm{hier}} =
\frac{1}{|\mathcal{T}|}
\sum_{(f,m,c)\in \mathcal{T}}
\Big[
1 - \operatorname{cos}(v_{\mathrm{sel}}^f, v_{\mathrm{sel}}^m) +
1 - \operatorname{cos}(v_{\mathrm{sel}}^m, v_{\mathrm{sel}}^c)
\Big].
\end{equation}
We regularize only adjacent levels, not all pairs, because the goal is not to collapse all temporal scopes into one representation. The goal is simply to prevent incompatible local evidence across nested supervision.

\subsection{Training Objective}
\label{sec:method_objective}

The global branch produces
\begin{equation}
v_{\mathrm{glob}} = \operatorname{norm}(P_v(A_v(E_v(c_{\mathrm{exp}})))).
\end{equation}
Both training branches share the same global text embedding $t$ and differ only in how the expanded clip is aggregated on the video side. Using $t$, we optimize two symmetric CLIP-style contrastive objectives:
\begin{equation}
\mathcal{L}_{\mathrm{glob}} =
\mathcal{L}_{v \rightarrow t}(v_{\mathrm{glob}}, t) +
\mathcal{L}_{t \rightarrow v}(t, v_{\mathrm{glob}}),
\end{equation}
\begin{equation}
\mathcal{L}_{\mathrm{sel}} =
\mathcal{L}_{v \rightarrow t}(v_{\mathrm{sel}}, t) +
\mathcal{L}_{t \rightarrow v}(t, v_{\mathrm{sel}}).
\end{equation}
The final objective is
\begin{equation}
\mathcal{L} =
\mathcal{L}_{\mathrm{glob}}
 + \lambda_{\mathrm{sel}} \mathcal{L}_{\mathrm{sel}}
 + \lambda_{\mathrm{hier}} \mathcal{L}_{\mathrm{hier}}.
\end{equation}

Each term has a distinct role. $\mathcal{L}_{\mathrm{glob}}$ preserves stable clip-level video-text alignment over the expanded training window. $\mathcal{L}_{\mathrm{sel}}$ teaches the model to recover local evidence within that window through text-guided reweighting. $\mathcal{L}_{\mathrm{hier}}$ prevents the selected branch from becoming inconsistent across neighboring semantic granularities.

\subsection{Inference}
\label{sec:method_inference}

At inference time, SurgAlign falls back to standard CLIP-style encoding:
\begin{equation}
v = \operatorname{encode\_image}(c), \qquad
t = \operatorname{encode\_text}(y),
\end{equation}
followed by cosine similarity in the shared embedding space. No expanded window, frame reweighting, or hierarchy-specific matching is required at test time. This separation is intentional: SurgAlign is a method for improving weak supervision during pretraining, not a query-conditioned inference system.


\section{Experiments}
\label{sec:experiments}

\subsection{Experimental Setup}
\label{sec:exp_data}

We pretrain on SurgLaVi-style weak video-text supervision, where each sample provides a video interval and a paired caption at fine, mid, or coarse granularity. Unless otherwise stated, SurgAlign uses the same backbone family, frame budget, optimizer, and contrastive objective as the matched baseline, and differs only in how weak temporal supervision is consumed. In the current primary experiment, we use a LemonFM-initialized visual encoder, a SurgicBERTa text encoder, 8 sampled frames, 2 GPUs with per-GPU batch size 128, AdamW with learning rate $10^{-4}$ and weight decay $0.02$, xpool-style frame selection, interval expansion ratio $1.5$, selection-loss weight $0.7$, hierarchical-consistency weight $0.1$, and mixed-level batches with a fine/mid/coarse ratio of $80/32/16$. During training, we sample multiple frames from an expanded temporal window, compute text-guided frame weights, and optimize both the global and selected branches. Test-time evaluation uses standard CLIP-style video/text encoding.

Our downstream evaluation follows the common zero-shot surgical recognition setup and covers phase, step, action, tool, and triplet prediction. This task diversity is important for our hypothesis: if temporal denoising truly helps recover local evidence from weakly aligned intervals, the benefit should be most visible on tasks that rely on fine-grained temporal cues rather than only broad workflow context.

The current results focus on the \textbf{primary setting}, which uses a surgical-domain backbone and serves as the main benchmark comparison against strong surgical VLP systems. This is the fairest setting for judging the method at its intended operating point. Complementary evidence analysis, ablation, and qualitative study are included to probe the mechanism beyond this primary comparison, and entries are marked as TBD where numerical results are omitted in the current draft.

\subsection{Main Benchmark Comparison}
\label{sec:exp_main}

The main benchmark comparison is designed around three questions: which tasks benefit most from SurgAlign, whether the method remains competitive with recent surgical VLP systems, and where the current design still underperforms stronger baselines. This framing is important because the method targets a specific bottleneck, temporal misalignment inside weak pairs, rather than attempting to solve every source of performance variation in surgical VLP. Our main claim in this setting is therefore not that SurgAlign must dominate every metric, but that under a surgical-domain backbone, misalignment-aware training should be most valuable on tasks that depend on localized temporal evidence.

\begin{table*}[t]
\centering
\caption{Main zero-shot comparison against representative surgical VLP baselines in the primary surgical-domain setting. Values are reported as Acc/F1 for phase, step, and action tasks, and mAP for triplet/tool tasks. Entries not yet available are marked as TBD.}
\label{tab:main_benchmark_comparison}
\resizebox{\textwidth}{!}{
\begin{tabular}{lllcccccccc}
\toprule
Recognition Task & Dataset & Metric & CLIP & MedSigLIP & SurgVLP & HecVL & PeskaVLP & SurgCLIP($\beta$) & SurgCLIP & Ours \\
\midrule
\multirow{12}{*}{Phase}
& \multirow{2}{*}{Cholec80} & Acc & TBD & TBD & TBD & TBD & TBD & 57.98 & 61.29 & 62.73 \\
&  & F1 & TBD & TBD & TBD & TBD & TBD & 39.42 & 50.53 & 43.46 \\
& \multirow{2}{*}{AutoLaparo} & Acc & TBD & TBD & TBD & TBD & TBD & 55.72 & 69.14 & 53.47 \\
&  & F1 & TBD & TBD & TBD & TBD & TBD & 45.95 & 56.37 & 48.48 \\
& \multirow{2}{*}{StrasBypass70} & Acc & TBD & TBD & TBD & TBD & TBD & 31.24 & 32.37 & 26.85 \\
&  & F1 & TBD & TBD & TBD & TBD & TBD & 26.05 & 30.78 & 23.87 \\
& \multirow{2}{*}{HeiChole} & Acc & TBD & TBD & TBD & TBD & TBD & 56.95 & 63.84 & 60.42 \\
&  & F1 & TBD & TBD & TBD & TBD & TBD & 44.00 & 55.15 & 47.42 \\
& \multirow{2}{*}{BernBypass70} & Acc & TBD & TBD & TBD & TBD & TBD & 18.30 & 23.90 & 16.56 \\
&  & F1 & TBD & TBD & TBD & TBD & TBD & 15.06 & 19.68 & 13.85 \\
& \multirow{2}{*}{GraSP} & Acc & TBD & TBD & TBD & TBD & TBD & 34.77 & 41.49 & 35.96 \\
&  & F1 & TBD & TBD & TBD & TBD & TBD & 27.98 & 34.94 & 32.29 \\
\midrule
\multirow{2}{*}{Step}
& \multirow{2}{*}{GraSP} & Acc & TBD & TBD & TBD & TBD & TBD & 14.15 & 26.28 & 19.58 \\
&  & F1 & TBD & TBD & TBD & TBD & TBD & 11.14 & 16.53 & 16.41 \\
\midrule
\multirow{2}{*}{Action}
& \multirow{2}{*}{SARRARP50} & Acc & TBD & TBD & TBD & TBD & TBD & 13.94 & 17.42 & 15.16 \\
&  & F1 & TBD & TBD & TBD & TBD & TBD & 7.62 & 7.76 & 7.67 \\
\midrule
Triplet & CholeT50 & mAP & TBD & TBD & TBD & TBD & TBD & 4.17 & 5.28 & 3.97 \\
\midrule
\multirow{3}{*}{Tool}
& Cholec80 & mAP & TBD & TBD & TBD & TBD & TBD & 36.77 & 40.80 & 33.39 \\
& HeiChole & mAP & TBD & TBD & TBD & TBD & TBD & 31.47 & 36.79 & 29.34 \\
& GraSP & mAP & TBD & TBD & TBD & TBD & TBD & 43.06 & 45.97 & 44.22 \\
\bottomrule
\end{tabular}
}
\end{table*}

Compared with SurgCLIP-$\beta$, SurgAlign improves multiple zero-shot benchmarks, including Cholec80 Phase ($+4.75$ Acc / $+4.04$ F1), HeiChole Phase ($+3.47$ Acc / $+3.42$ F1), GraSP Phase ($+1.19$ Acc / $+4.31$ F1), GraSP Step ($+5.43$ Acc / $+5.27$ F1), SARRARP50 Action ($+1.22$ Acc / $+0.05$ F1), and GraSP Tool ($+1.16$ mAP). The clearest gain appears on GraSP Step, which is consistent with the hypothesis that misalignment-aware local reweighting is most useful for tasks requiring localized temporal evidence. At the same time, SurgAlign does not uniformly surpass the stronger full SurgCLIP baseline. The current result should therefore be interpreted as encouraging evidence for the value of temporal denoising, rather than as a claim of consistent dominance over stronger end-to-end surgical VLP baselines.

\subsection{Expanded-Window Evidence Analysis}
\label{sec:exp_evidence}

The second experiment asks whether the expanded temporal neighborhood is actually used by the selection mechanism, rather than serving as a formal augmentation. After training, we partition sampled frames into those inside the original weak interval and those in the expanded-only region, then aggregate four statistics: outside-frame fraction, outside mass, outside-evidence enrichment relative to the uniform frame fraction, and the rate at which the top-weighted frame falls outside the original interval. Under weak supervision, this analysis is not meant to certify outside frames as ground-truth positives. It tests the narrower claim that SurgAlign assigns meaningful and enriched evidence mass beyond the noisy annotated boundary.

\begin{table}[t]
\centering
\caption{Expanded-window evidence analysis. For each supervision level, we report outside-frame fraction, outside mass, outside-evidence enrichment, and top-1 outside rate. Entries are marked as TBD where numerical results are omitted in the current draft.}
\label{tab:expanded_window_evidence}
\resizebox{\textwidth}{!}{
\begin{tabular}{lcccc}
\toprule
Level & Outside Frame Fraction & Outside Mass & Outside Enrichment & Top-1 Outside Rate \\
\midrule
Fine & TBD & TBD & TBD & TBD \\
Mid & TBD & TBD & TBD & TBD \\
Coarse & TBD & TBD & TBD & TBD \\
\bottomrule
\end{tabular}
}
\end{table}

\subsection{Ablation Study}
\label{sec:exp_ablation}

The ablation study is conducted in the primary setting so that the only changing factor is the training rule. The most important question here is not whether SurgAlign beats a large set of literature baselines, but which parts of the proposed pipeline are necessary. We therefore ablate the method progressively: expanded temporal evidence, text-guided local weighting, the selected-view auxiliary loss, and hierarchical consistency.

\begin{table}[t]
\centering
\caption{Ablation study in the primary setting. The table isolates the contributions of expanded temporal evidence, text-guided local reweighting, the selected auxiliary branch, and hierarchical consistency. Entries are marked as TBD where numerical results are omitted in the current draft.}
\label{tab:ablation_multitask}
\resizebox{\textwidth}{!}{
\begin{tabular}{lllccccc}
\toprule
Recognition Task & Dataset & Metric & Baseline & + Expand & + Reweight & + Sel. Loss & Full \\
\midrule
\multirow{12}{*}{Phase}
& \multirow{2}{*}{Cholec80} & Acc & TBD & TBD & TBD & TBD & TBD \\
&  & F1 & TBD & TBD & TBD & TBD & TBD \\
& \multirow{2}{*}{AutoLaparo} & Acc & TBD & TBD & TBD & TBD & TBD \\
&  & F1 & TBD & TBD & TBD & TBD & TBD \\
& \multirow{2}{*}{StrasBypass70} & Acc & TBD & TBD & TBD & TBD & TBD \\
&  & F1 & TBD & TBD & TBD & TBD & TBD \\
& \multirow{2}{*}{Heichole} & Acc & TBD & TBD & TBD & TBD & TBD \\
&  & F1 & TBD & TBD & TBD & TBD & TBD \\
& \multirow{2}{*}{BernBypass70} & Acc & TBD & TBD & TBD & TBD & TBD \\
&  & F1 & TBD & TBD & TBD & TBD & TBD \\
& \multirow{2}{*}{GraSP} & Acc & TBD & TBD & TBD & TBD & TBD \\
&  & F1 & TBD & TBD & TBD & TBD & TBD \\
\midrule
\multirow{2}{*}{Step}
& \multirow{2}{*}{GraSP} & Acc & TBD & TBD & TBD & TBD & TBD \\
&  & F1 & TBD & TBD & TBD & TBD & TBD \\
\midrule
\multirow{2}{*}{Action}
& \multirow{2}{*}{SARRARP50} & Acc & TBD & TBD & TBD & TBD & TBD \\
&  & F1 & TBD & TBD & TBD & TBD & TBD \\
\midrule
Triplet & CholecT50 & mAP & TBD & TBD & TBD & TBD & TBD \\
\midrule
\multirow{3}{*}{Tool}
& Cholec80 & mAP & TBD & TBD & TBD & TBD & TBD \\
& Heichole & mAP & TBD & TBD & TBD & TBD & TBD \\
& GraSP & mAP & TBD & TBD & TBD & TBD & TBD \\
\bottomrule
\end{tabular}
}
\end{table}

\subsection{Qualitative Study}
\label{sec:exp_qualitative}

We also include representative qualitative case studies from the expanded window to examine whether the learned weights focus on visually plausible local evidence.

\paragraph{Frame-weight visualization.}
For each case, we visualize sampled frames, their learned weights, the original weak interval, and the paired caption. We include both successful and challenging examples, including cases where the highest-weight frame lies outside the original weak interval.


\section{Discussion and Limitations}
\label{sec:discussion}

SurgAlign is best understood as a supervision-refinement method. It targets one specific failure mode of weak video-language pretraining: a video-text pair can be semantically valid even when only a small subset of frames provides the visual evidence that justifies it. For that reason, the method is intentionally training-only. We do not introduce a new test-time retrieval head or a query-conditioned inference pipeline; the claim is simply that weak temporal positives should be handled more carefully during pretraining.

The evaluation protocol is designed around the first limitation. SurgAlign is most appropriate when the paired text remains semantically informative, the temporal offset is local rather than arbitrarily long-range, and the decisive evidence is sparse enough that uniform pooling becomes harmful. The method is therefore not a substitute for stronger backbones, richer negatives, or better corpus design.

The second limitation is empirical scope. Our main claim is tied to internal matched comparisons, with SurgCLIP-$\beta$ and full SurgCLIP treated as reported strong references rather than as matched baselines. This distinction matters because temporal misalignment is real but not exclusive: tasks dominated by broader workflow context, harder negative discrimination, or label ambiguity can require additional modeling beyond local frame weighting.

A third limitation is optimization sensitivity. In practice, the final behavior can depend on the expansion ratio, the auxiliary loss weight, and the training trajectory. To avoid converting this sensitivity into hidden test-set checkpoint selection, we fix the official checkpoint at epoch 50 and report multi-seed aggregates. Additional checkpoint trajectories are used only as diagnostic analyses.

A final limitation is evaluation breadth. We validate SurgAlign on surgical-domain benchmarks, where temporal misalignment is practically important and visually subtle, but we do not yet show transfer to broader narrated video datasets. The broader claim of the paper is therefore not that SurgAlign solves weak video-language alignment in general, but that treating weak intervals as noisy temporal anchors is a useful training principle in a challenging real-world regime.


\section{Conclusion}
\label{sec:conclusion}

We introduced SurgAlign, a misalignment-aware framework for surgical video-language pretraining. The central idea is that weakly paired surgical intervals should be treated as noisy temporal anchors rather than uniformly reliable positive clips. SurgAlign operationalizes this idea through expanded temporal evidence, text-guided local reweighting, and lightweight hierarchical consistency, while preserving standard CLIP-style inference.

This framing makes the contribution precise. We do not claim a new surgical backbone or a new inference protocol. Instead, we show how to improve surgical VLP by modeling temporal misalignment inside weak video-text pairs during training. In the primary surgical-domain setting, the method improves several zero-shot benchmarks while remaining compatible with standard CLIP-style inference.

More broadly, the paper argues that future progress in weakly supervised surgical VLP should not be driven only by larger corpora or stronger encoders. It should also come from better modeling of how weak temporal supervision is consumed. SurgAlign is one step in that direction.


\bibliographystyle{IEEEtran}
\bibliography{related_work_refs}

\end{document}
